% ---------------------------------------------------------------------------
% Environment library overview
% ---------------------------------------------------------------------------
\begingroup
% Line breaking for paragraphs with long, unbreakable identifiers (local to
% this chapter): allow looser lines instead of overfull ones.
\setlength{\emergencystretch}{3em}\tolerance=1000
\chapter{Environment Library Overview}\label{chap:envlib}

The \pkg{env\_lib} package provides multi-agent environments for research on
networked control, cooperative decision making and multi-robot estimation. All
environments implement the Gymnasium API, are registered with Gymnasium when
\pkg{env\_lib} is imported, and share one registry, one rendering and
recording layer and one set of graph utilities. This chapter describes what
the environments have in common; the following chapters document each
environment in detail:

\begin{itemize}
  \item Kuramoto oscillator networks, NumPy and batched PyTorch backends
        (Chapter~\ref{chap:kuramoto});
  \item LineMsg and WirelessComm, two discrete networked communication
        benchmarks (Chapter~\ref{chap:networked});
  \item Pistonball, a cooperative physics game (Chapter~\ref{chap:pistonball});
  \item Consensus and formation control on communication graphs
        (Chapter~\ref{chap:consensus});
  \item PowerGrid, frequency control of a networked power system
        (Chapter~\ref{chap:power_grid});
  \item Platoon, cooperative adaptive cruise control of a vehicle platoon
        (Chapter~\ref{chap:platoon});
  \item AJLATT, active joint localisation and target tracking with a robot
        team, including its map tools (Chapter~\ref{chap:ajlatt}).
\end{itemize}

Chapter~\ref{chap:workflow} then describes the layer that connects the
environments to experiments: native vector environments, the catalogue and
the \code{env-lib} command line, adapters, baseline controllers and
evaluation.

% ---------------------------------------------------------------------------
\section{Environments at a glance}\label{sec:envlib-comparison}

Table~\ref{tab:envlib-comparison} summarises the environments in their default
configuration. $N$ is the number of oscillators, $M$ the number of coupling
edges and $n$ the number of agents. The NumPy Kuramoto, Consensus, Formation,
PowerGrid and Platoon environments have native vector implementations
(Section~\ref{sec:envlib-vector}).

\begin{table}[htbp]
  \centering\footnotesize
  \setlength{\tabcolsep}{4pt}
  \caption{Environments in their default configuration.}
  \label{tab:envlib-comparison}
  \begin{tabularx}{\textwidth}{@{}>{\raggedright\arraybackslash}p{2.2cm}>{\raggedright\arraybackslash}p{1.9cm}>{\raggedright\arraybackslash}X>{\raggedright\arraybackslash}p{2.0cm}>{\raggedright\arraybackslash}p{2.6cm}>{\raggedright\arraybackslash}p{1.5cm}@{}}
    \toprule
    Environment & Agents & Action space & Observation & Reward & Requires \\
    \midrule
    Kuramoto (NumPy) & 1 controller, $N = 10$ oscillators &
      continuous \code{Box((N+M,))}, $N + M = 55$ & \code{(75,)}, i.e.\ $3N + M$ &
      scalar synchronisation measure, bonus on synchronisation & none \\
    Kuramoto (PyTorch) & \code{n\_agents} parallel systems (1) &
      as NumPy; batched \code{(n\_agents, N+M)} accepted & \code{(75,)} of system 0 &
      per system in \code{info} & \pkg{torch} \\
    LineMsg & $n = 10$ & \code{Discrete(2**n)} or \code{MultiBinary(n)} &
      \code{(10, 3)} & sum of per-agent rewards & none \\
    WirelessComm & $n = 36$ ($6 \times 6$) & \code{MultiDiscrete([5]*n)} &
      \code{(36, 18)} & packets delivered (sum) & none \\
    Pistonball & $n = 20$ pistons & \code{Box(-1, 1, (n,))} or \code{MultiDiscrete([3]*n)} &
      \code{(20, 7)} & sum of local rewards & \pkg{pymunk}, \pkg{pygame} \\
    Consensus / Formation & $n = 8$ & \code{Box(-1, 1, (n, 2))} &
      \code{(8, 16)} & mean of local rewards, success bonus & none \\
    PowerGrid & $n = 16$ buses & \code{Box(-0.1, 0.1, (n,))}: injection in per unit &
      \code{(16, 10)} & sum of local rewards, trip penalty & none \\
    Platoon & $n = 8$ followers & \code{Box(-6, 3, (n,))}: acceleration command in m/s$^2$ &
      \code{(8, 13)} & sum of local rewards, collision and break-up penalties & none \\
    AJLATT & $n = 4$ robots & \code{Box((n, 2))}: $v \in [0, 2]$, $\omega \in [-\pi/4, \pi/4]$ &
      \code{(4, 24)} & per-robot array & none \\
    \bottomrule
  \end{tabularx}
\end{table}

% ---------------------------------------------------------------------------
\section{Package structure}\label{sec:envlib-structure}

The package lives in \code{src/env\_lib} of the repository:

\begin{textcode}
src/env_lib/
  __init__.py           public names, lazily imported classes and functions
  registration.py       ids and metadata (EnvSpec), make(), make_vec(), get_spec()
  catalog.py            catalog(), describe(), EnvInfo
  baselines.py          baseline controllers, baseline_policy(), list_baselines()
  __main__.py           the env-lib command line (python -m env_lib)
  errors.py             ResetNeededError
  wrappers/             FlattenJointSpaces, TeamReward, ParallelEnvAdapter
  utils/                themes, renderers, recording, graphs, BatchedVectorEnv,
                        evaluate() and rollout()
  kos_env/              Kuramoto oscillators (NumPy, PyTorch, native vector)
  consensus_env/        consensus and formation control (single, native vector)
  power_grid_env/       PowerGrid: frequency control (single, native vector)
  platoon_env/          Platoon: cooperative adaptive cruise control
  ajlatt_env/           AJLATT: multi-robot localisation and target tracking
    maps/               occupancy maps, map builder and map plotting
  pistonball_env/       Pistonball: cooperative physics game (pymunk, pygame)
  linemsg_env/          LineMsg: message relay along a line of agents
  wireless_comm_env/    WirelessComm: random access to shared access points
\end{textcode}

Every environment subpackage contains the environment module and a
\code{rendering} module. The renderer is imported on the first call to
\code{render()}, so an environment that is never rendered does not create any
figure or window.

\paragraph{Lazy imports.} \code{import env\_lib} imports neither PyTorch nor
pygame, pymunk or tkinter. The environment classes of
Table~\ref{tab:envlib-classes} are attributes of the package that are imported
on first access, for example \code{env\_lib.ConsensusEnv}; the same holds for
\code{env\_lib.baseline\_policy}, \code{env\_lib.evaluate},
\code{env\_lib.rollout}, \code{env\_lib.describe} and \code{env\_lib.EnvInfo}
and for the subpackages \pkg{env\_lib.wrappers}, \pkg{env\_lib.baselines} and
\pkg{env\_lib.catalog}. An optional dependency is therefore only needed when
the environment that uses it is created (Table~\ref{tab:envlib-optional}).

\begin{table}[htbp]
  \centering\small
  \caption{Environment classes exported by \pkg{env\_lib}.}
  \label{tab:envlib-classes}
  \begin{tabularx}{\textwidth}{@{}>{\raggedright\arraybackslash}p{5.3cm}>{\raggedright\arraybackslash}X>{\raggedright\arraybackslash}p{0.7cm}@{}}
    \toprule
    Class & Module & Ch. \\
    \midrule
    \code{KuramotoOscillatorEnv} & \code{env\_lib.kos\_env.kuramoto\_env} & \ref{chap:kuramoto} \\
    \code{KuramotoOscillatorEnvTorch} & \code{env\_lib.kos\_env.kuramoto\_env\_torch} & \ref{chap:kuramoto} \\
    \code{KuramotoOscillatorVectorEnv} & \code{env\_lib.kos\_env.vector} & \ref{chap:workflow} \\
    \code{LineMsgEnv} & \code{env\_lib.linemsg\_env.linemsg\_env} & \ref{chap:networked} \\
    \code{WirelessCommEnv} & \code{env\_lib.wireless\_comm\_env.wireless\_comm\_env} & \ref{chap:networked} \\
    \code{PistonballEnv} & \code{env\_lib.pistonball\_env.pistonball\_env} & \ref{chap:pistonball} \\
    \code{ConsensusEnv} & \code{env\_lib.consensus\_env.consensus\_env} & \ref{chap:consensus} \\
    \code{ConsensusVectorEnv} & \code{env\_lib.consensus\_env.vector} & \ref{chap:workflow} \\
    \code{PowerGridEnv}, \code{PowerGridVectorEnv} & \code{env\_lib.power\_grid\_env.power\_grid\_env} & \ref{chap:power_grid} \\
    \code{PlatoonEnv}, \code{PlatoonVectorEnv} & \code{env\_lib.platoon\_env.platoon\_env} & \ref{chap:platoon} \\
    \code{AJLATTEnv}, \code{AJLATTConfig} & \code{env\_lib.ajlatt\_env} & \ref{chap:ajlatt} \\
    \bottomrule
  \end{tabularx}
\end{table}

\begin{table}[htbp]
  \centering\small
  \caption{Optional dependencies of the environments. The extras are
    installed with \code{pip install -e ".[extra]"} from the repository root.}
  \label{tab:envlib-optional}
  \begin{tabularx}{\textwidth}{@{}ll>{\raggedright\arraybackslash}X@{}}
    \toprule
    Extra & Packages & Needed for \\
    \midrule
    \code{pistonball} & \pkg{pymunk}, \pkg{pygame} & Pistonball: \pkg{pymunk} to create the
      environment, \pkg{pygame} to render it and for keyboard control \\
    \code{torch} & \pkg{torch} & the PyTorch backend of the Kuramoto environments \\
    \code{video} & \pkg{imageio}, \pkg{imageio-ffmpeg} & \code{save\_animation} to \code{.mp4} and
      other video formats (GIF export needs only Pillow) \\
    \bottomrule
  \end{tabularx}
\end{table}

PettingZoo is not an extra: the parallel adapter of
Section~\ref{sec:workflow-wrappers} works without it and subclasses
PettingZoo's \code{ParallelEnv} when it is installed
(Section~\ref{sec:install-pettingzoo}).

% ---------------------------------------------------------------------------
\section{Creating environments}\label{sec:envlib-make}

There are three equivalent ways to create an environment:

\begin{pycode}
import gymnasium as gym
import env_lib

env = env_lib.make("WirelessComm-v0")               # registered id
env = gym.make("WirelessComm-v0")                   # after importing env_lib
env = env_lib.WirelessCommEnv(grid_x=6, grid_y=6)   # direct construction

print(len(env_lib.list_envs()))                     # 18 registered ids
\end{pycode}

The function \code{env\_lib.make} registers the environments if necessary and
calls \code{gymnasium.make} with the id and the keyword arguments, which are
passed to the constructor and override the configuration of the id, for example
\code{env\_lib.make("LineMsg-v0", num\_agents=20, max\_iter=100)}.
\code{gymnasium.make} returns the environment inside Gymnasium's standard
wrappers (order enforcement and, except for \envid{AJLATT-v0}, the passive
environment checker). Environment-specific methods and attributes such as
\code{laplacian\_policy()}, \code{adjacency} or \code{observation\_layout} are
reached through \code{env.unwrapped}.

\paragraph{Episode length.} \code{env\_lib.make(env\_id, max\_episode\_steps=N)}
sets the environment's own limit argument (\code{max\_steps},
\code{max\_iter}, \code{max\_cycles} or \code{max\_episode\_steps}, recorded in
\code{EnvSpec.limit\_kwarg}) to \code{N} instead of adding a \code{TimeLimit}
wrapper, so there is always exactly one limit. Passing both
\code{max\_episode\_steps} and a different value of the limit argument raises
\code{ValueError}. \code{gymnasium.make} itself does not translate the
argument (Section~\ref{sec:envlib-episodes}).

\paragraph{Vector environments.} \code{env\_lib.make\_vec(env\_id, num\_envs)}
creates \code{num\_envs} copies as one Gymnasium vector environment: the
environment's native batched implementation when it has one, and
\code{gymnasium.vector.SyncVectorEnv} otherwise
(Section~\ref{sec:envlib-vector}). It takes the same keyword arguments as
\code{make}, including \code{max\_episode\_steps}.

\begin{pycode}
import env_lib

env = env_lib.make("Platoon-v0", max_episode_steps=300)
print(env.unwrapped.max_steps, env.spec.max_episode_steps)   # 300 None

envs = env_lib.make_vec("Platoon-v0", num_envs=64)      # native batch
print(type(envs).__name__)                              # PlatoonVectorEnv
print(envs.observation_space.shape)                     # (64, 8, 13)
envs = env_lib.make_vec("LineMsg-v0", num_envs=4)       # no native version
print(type(envs).__name__)                              # SyncVectorEnv
\end{pycode}

% ---------------------------------------------------------------------------
\section{Registry and metadata}\label{sec:envlib-registry}

Every registered id is described by an \code{EnvSpec} record (class
\code{env\_lib.registration.EnvSpec}) in \code{env\_lib.registration.ENV\_SPECS};
\code{env\_lib.get\_spec(env\_id)} returns the record of one id and raises
\code{KeyError} for an unknown id. \code{EnvSpec} is a frozen dataclass,
distinct from Gymnasium's \code{gymnasium.envs.registration.EnvSpec} (which
\code{env.spec} returns); its fields are listed in
Table~\ref{tab:envlib-spec}. The catalogue (\code{env\_lib.catalog()},
Section~\ref{sec:workflow-catalog}) and the \code{env-lib} command line are
built on these records, so an environment can be selected by its properties
instead of by name.

\begin{table}[htbp]
  \centering\small
  \caption{Fields of \code{env\_lib.registration.EnvSpec}.}
  \label{tab:envlib-spec}
  \begin{tabularx}{\textwidth}{@{}>{\raggedright\arraybackslash}p{3.6cm}>{\raggedright\arraybackslash}X@{}}
    \toprule
    Field (default) & Content \\
    \midrule
    \code{id} & registered Gymnasium id \\
    \code{entry\_point} & \code{"module:Class"} of the single environment \\
    \code{kwargs} (\code{\{\}}) & constructor arguments of this configuration \\
    \code{description} (\code{""}) & one-line summary \\
    \code{family} (\code{""}) & environment family: the subpackage name without
      \code{\_env} (\code{"kuramoto"}, \code{"linemsg"}, \code{"wireless\_comm"},
      \code{"pistonball"}, \code{"consensus"}, \code{"power\_grid"},
      \code{"platoon"}, \code{"ajlatt"}) \\
    \code{observation\_type}, \code{action\_type} (\code{"continuous"}) &
      \code{"continuous"}, \code{"discrete"} or, for actions that can be either
      (Pistonball), \code{"continuous|discrete"} \\
    \code{requires} (\code{None}) & optional-dependency extra needed by the
      environment (\code{"torch"} or \code{"pistonball"}) \\
    \code{vector\_entry\_point} (\code{None}) & \code{"module:Class"} of the native
      batched implementation used by \code{make\_vec} \\
    \code{limit\_kwarg} (\code{"max\_steps"}) & name of the constructor argument that
      sets the episode length; \code{make} and \code{make\_vec} translate
      \code{max\_episode\_steps} into it \\
    \code{per\_agent\_rewards} (\code{False}) & rewards and termination flags are
      per-agent arrays (AJLATT); \code{make\_vec} then wraps Sync and Async copies
      in \code{TeamReward} \\
    \code{disable\_env\_checker} (\code{False}) & skip Gymnasium's passive checker
      (\code{True} for AJLATT, whose rewards are per-robot arrays) \\
    \bottomrule
  \end{tabularx}
\end{table}

\begin{pycode}
import env_lib

spec = env_lib.get_spec("PowerGrid-v0")
print(spec.family, spec.action_type)   # power_grid continuous
print(spec.limit_kwarg)                # max_steps
print(spec.vector_entry_point)
# env_lib.power_grid_env.power_grid_env:PowerGridVectorEnv

native = env_lib.catalog(native_vector=True, inspect=False).ids()
print(len(native))                     # 8 ids with a native batch
\end{pycode}

% ---------------------------------------------------------------------------
\section{Common API conventions}\label{sec:envlib-api}

\subsection{Gymnasium interface}

Every environment implements
\begin{itemize}
  \item \code{reset(*, seed=None, options=None)} returning
        \code{(observation, info)};
  \item \code{step(action)} returning
        \code{(observation, reward, terminated, truncated, info)};
  \item \code{render()} and \code{close()}; \code{close()} may be called
        several times.
\end{itemize}
Observations are \code{float32} NumPy arrays contained in
\code{observation\_space}. Calling \code{step()} before the first
\code{reset()} raises \code{env\_lib.ResetNeededError}, which derives from
both \code{RuntimeError} and \code{gymnasium.error.ResetNeeded}, so either can
be caught; the native vector environments raise it as well, also when the
first \code{reset()} resets only some copies. Invalid
actions (wrong shape, values outside a discrete range, NaN or infinite values)
raise \code{ValueError}, and continuous actions are clipped to the action
bounds. AJLATT is the exception: it only checks the action shape, and clipping
is opt-in (\code{clip\_actions=True}, Section~\ref{sec:ajlatt-spaces}).

\subsection{Joint multi-agent spaces}\label{sec:envlib-joint}

Each environment exposes the whole team through a single Gymnasium interface.
The action space covers all agents at once (for example
\code{Box(-1, 1, (8, 2))} for eight planar agents,
\code{Box(-0.1, 0.1, (16,))} for the 16 buses of a power grid, or
\code{MultiDiscrete([5] * 36)} for 36 wireless agents), and observations have
one row per agent. A multi-agent algorithm splits the joint arrays along the
first axis; a centralised algorithm uses them directly. The wrappers of
Section~\ref{sec:workflow-wrappers} provide both views ready-made: a flat
vector for single-agent libraries and per-agent dictionaries for the
PettingZoo Parallel API.

The scalar \code{reward} returned by \code{step()} is the team reward. The
per-agent rewards are always available as \code{info["agent\_rewards"]}, a
\code{float64} array of shape \code{(n\_agents,)}, which every \code{step()}
returns. How the team reward is formed from the per-agent rewards depends on the
environment (Table~\ref{tab:envlib-comparison}): it is their sum for LineMsg,
WirelessComm, Pistonball, PowerGrid, Platoon and AJLATT and their mean (plus a
success bonus) for Consensus. The Kuramoto environments control one or several
complete oscillator networks, so their ``agents'' are systems: the NumPy
backend reports \code{n\_agents} copies of its reward and the PyTorch backend
one reward per simulated system. AJLATT returns the per-robot rewards and
per-robot collision flags directly as arrays; \code{TeamReward}
(Section~\ref{sec:workflow-wrappers}) and the older \code{TeamRewardWrapper}
(Section~\ref{sec:ajlatt-wrapper}) convert them to a scalar reward and a scalar
\code{terminated} flag.

\subsection{Info dictionaries}\label{sec:envlib-info}

The info dictionary returned by \code{reset()} contains
\code{"agent\_rewards"} (all zeros) for LineMsg, WirelessComm, Consensus,
PowerGrid, Platoon and AJLATT. The Kuramoto and Pistonball environments do not
include the key in their reset info. Table~\ref{tab:envlib-info} lists the
other keys of the step info; the environment chapters define them.

\begin{table}[htbp]
  \centering\small
  \caption{Keys of the info dictionary returned by \code{step()} besides
    \code{"agent\_rewards"}.}
  \label{tab:envlib-info}
  \begin{tabularx}{\textwidth}{@{}>{\raggedright\arraybackslash}p{2.4cm}>{\raggedright\arraybackslash}X>{\raggedright\arraybackslash}p{1.5cm}@{}}
    \toprule
    Environment & Keys & Details \\
    \midrule
    Kuramoto & \code{order\_parameter}, \code{phase\_coherence}, \code{phases},
      \code{natural\_frequencies}, \code{coupling\_matrix}, \code{dphases\_dt},
      \code{step\_count}, \code{device}, \code{n\_agents} &
      Sec.~\ref{sec:kuramoto-info} \\
    LineMsg & none & Sec.~\ref{sec:linemsg-params} \\
    WirelessComm & \code{outcomes}, \code{ap\_load} & Sec.~\ref{sec:wireless-params} \\
    Pistonball & \code{local\_rewards}, \code{total\_reward} & Sec.~\ref{sec:pistonball-params} \\
    Consensus & \code{error}, \code{algebraic\_connectivity}, \code{adjacency},
      \code{success}, \code{step} & Sec.~\ref{sec:consensus-info} \\
    PowerGrid & \code{omega}, \code{rocof}, \code{max\_abs\_omega},
      \code{frequency\_nadir}, \code{peak\_abs\_omega}, \code{coi\_omega},
      \code{frequency\_spread}, \code{control\_effort}, \code{order\_parameter},
      \code{total\_disturbance}, \code{tripped}, \code{step}, \code{time} &
      Sec.~\ref{sec:power-grid-info} \\
    Platoon & \code{spacing\_errors}, \code{gaps}, \code{relative\_speeds},
      \code{speeds}, \code{accelerations}, \code{commands},
      \code{peak\_spacing\_errors}, \code{error\_amplification},
      \code{leader\_speed}, \code{leader\_acceleration}, \code{min\_gap},
      \code{collision}, \code{breakup}, \code{step}, \code{time} &
      Sec.~\ref{sec:platoon-info} \\
    AJLATT & \code{team\_reward}, \code{collisions}, \code{episode\_collisions},
      \code{target\_observed}, \code{communication}, \code{robot\_cov\_trace},
      \code{target\_cov\_trace}, \code{step}, \code{numerical\_error} &
      Sec.~\ref{sec:ajlatt-info} \\
    \bottomrule
  \end{tabularx}
\end{table}

In a vector environment every info value gains a leading \code{num\_envs}
dimension and comes with Gymnasium's boolean mask \code{"\_key"}
(Section~\ref{sec:workflow-vector}).

\subsection{Termination and truncation}\label{sec:envlib-episodes}

\code{terminated} signals a task-specific terminal state; \code{truncated}
signals the end of the time budget. Every environment enforces its own episode
limit through a constructor argument (Table~\ref{tab:envlib-limits}). The
registrations do not add a \code{TimeLimit} wrapper (\code{env.spec.max\_episode\_steps}
is \code{None}), so the limit argument is the only limit. To change the episode
length, pass the constructor argument, for example
\code{env\_lib.make("Consensus-v0", max\_steps=400)}, or pass
\code{max\_episode\_steps} to \code{env\_lib.make} or \code{env\_lib.make\_vec},
which set the same argument. Only \code{gymnasium.make(...,
max\_episode\_steps=N)} called directly adds an independent \code{TimeLimit}
wrapper on top of the environment's own limit; the shorter of the two then
applies.

\begin{table}[htbp]
  \centering\small
  \caption{Episode limits and termination conditions.}
  \label{tab:envlib-limits}
  \begin{tabularx}{\textwidth}{@{}lll>{\raggedright\arraybackslash}X@{}}
    \toprule
    Environment & Limit argument & Default & \code{terminated} is set when \\
    \midrule
    Kuramoto & \code{max\_steps} & 1000 & the order parameter exceeds \code{sync\_threshold} \\
    LineMsg & \code{max\_iter} & 50 & never (always \code{False}) \\
    WirelessComm & \code{max\_iter} & 50 & never (always \code{False}) \\
    Pistonball & \code{max\_cycles} & 125 & the ball reaches the left wall \\
    Consensus / Formation & \code{max\_steps} & 200 & the task error drops below \code{tolerance} \\
    PowerGrid & \code{max\_steps} & 200 & a frequency deviation exceeds the trip limit,
      $\max_i |\omega_i| > 2\pi$\,\code{frequency\_limit} \\
    Platoon & \code{max\_steps} & 600 & a follower collides with its predecessor or its
      gap exceeds \code{breakup\_distance} \\
    AJLATT & \code{max\_episode\_steps} & 120 & per robot: an obstacle is closer than
      \code{obstacle\_collision\_distance} (array of flags) \\
    \bottomrule
  \end{tabularx}
\end{table}

\subsection{Seeding and reproducibility}\label{sec:envlib-seeding}

\begin{itemize}
  \item \code{reset(seed=s)} seeds the environment's generator
        \code{self.np\_random}. All randomness of an episode (initial state,
        stochastic transitions, noise, load disturbances, leader profiles)
        is drawn from this generator, and \code{reset()} without a seed
        continues its stream. With the same seed, the same actions and the
        same library versions an episode is reproduced exactly.
  \item No environment reads or modifies the global NumPy or PyTorch random
        state. The PyTorch Kuramoto backend owns a \code{torch.Generator} that
        is seeded from \code{self.np\_random} at every reset.
  \item \code{env.action\_space.sample()} uses the space's own generator; seed
        it separately with \code{env.action\_space.seed(s)}.
  \item Some structure is fixed at construction and has its own seed:
        \code{topology\_seed} (random Kuramoto topology, default 42),
        \code{graph\_seed} (Erd\H{o}s--R\'enyi graph of the consensus
        environment), \code{network\_seed} (PowerGrid network and physical
        parameters, default 0), \code{vehicle\_seed} (Platoon vehicle lags and
        lengths, default 0) and \code{AJLATTConfig.seed} (seed of the first
        AJLATT reset when none is given). With \code{randomize\_network=True}
        (PowerGrid) or \code{randomize\_vehicles=True} (Platoon) this
        structure is instead resampled from \code{self.np\_random} at every
        reset.
  \item A native vector environment draws all copies from one generator, so
        \code{reset(seed=s)} makes the batch reproducible, but copy $i$ does
        not reproduce a single environment seeded with $s + i$
        (Section~\ref{sec:workflow-vector}).
  \item The legacy \code{env.seed()} methods of LineMsg, WirelessComm and
        Pistonball still work but emit a \code{DeprecationWarning}; use
        \code{reset(seed=...)} instead.
\end{itemize}

\subsection{Reset options}\label{sec:envlib-reset-options}

\code{reset(options=...)} can override parts of the sampled initial state
(Table~\ref{tab:envlib-options}). The values of the listed keys are validated,
and a wrong shape or a non-finite entry raises \code{ValueError}.
\emph{Unknown keys are ignored with a} \code{UserWarning} that names them and
the supported keys; they do not raise, so generic tools that pass options of
their own (such as PettingZoo's \code{parallel\_api\_test}) work with every
environment. LineMsg, WirelessComm, Pistonball and AJLATT have no reset
options and ignore \code{options} altogether. The native vector environments
accept the same keys, as arrays for one copy (applied to every reset copy) or
with a leading \code{num\_envs} dimension, and in addition \code{"reset\_mask"}
(Section~\ref{sec:workflow-vector}).

\begin{table}[htbp]
  \centering\small
  \caption{Reset options.}
  \label{tab:envlib-options}
  \begin{tabularx}{\textwidth}{@{}>{\raggedright\arraybackslash}p{3.0cm}>{\raggedright\arraybackslash}X@{}}
    \toprule
    Environment & Keys \\
    \midrule
    Kuramoto & \code{"phases"}, \code{"natural\_frequencies"}, \code{"coupling\_strengths"} (dynamic coupling only) \\
    Consensus / Formation & \code{"positions"}, \code{"velocities"}, each \code{(n, 2)} \\
    PowerGrid & \code{"theta"}, \code{"omega"}, \code{"disturbance"}, each \code{(n,)} \\
    Platoon & \code{"positions"} and \code{"velocities"} (given together), \code{"accelerations"},
      \code{"time\_constants"}, \code{"lengths"}, each \code{(n + 1,)} with the leader first;
      \code{"leader\_command"} \code{(max\_steps,)} \\
    others & none; \code{options} is ignored \\
    \bottomrule
  \end{tabularx}
\end{table}

\begin{pycode}
import warnings
import numpy as np
import env_lib

env = env_lib.make("PowerGrid-v0")
omega = np.full(16, 0.1)                      # 0.1 rad/s at every bus
with warnings.catch_warnings(record=True) as caught:
    warnings.simplefilter("always")
    obs, info = env.reset(seed=0, options={"omega": omega, "tag": "run-1"})
print(info["omega"][:3])                      # [0.1 0.1 0.1]
print(caught[0].category.__name__)            # UserWarning ("tag" is ignored)
\end{pycode}

% ---------------------------------------------------------------------------
\section{Vectorised simulation and the workflow layer}\label{sec:envlib-vector}

\code{env\_lib.make\_vec} returns a Gymnasium vector environment in which
observations, rewards and flags have a leading \code{num\_envs} dimension. For
the NumPy Kuramoto ids, \envid{Consensus-v0}, \envid{Formation-v0},
\envid{PowerGrid-v0} and \envid{Platoon-v0} it is a native implementation
(\code{KuramotoOscillatorVectorEnv}, \code{ConsensusVectorEnv},
\code{PowerGridVectorEnv}, \code{PlatoonVectorEnv}) that advances all copies
with one set of array operations and shares its dynamics code with the single
environment. The native classes derive from
\code{env\_lib.utils.BatchedVectorEnv}, which implements seeding, action
validation, the three autoreset modes and the info masks on top of three
array hooks. All other ids are wrapped in \code{gymnasium.vector.SyncVectorEnv}
(\code{vectorization\_mode="async"} gives an \code{AsyncVectorEnv}).

Chapter~\ref{chap:workflow} documents this layer in full: vector environments,
autoreset and seeding (Section~\ref{sec:workflow-vector}), the catalogue and
the command line (Sections~\ref{sec:workflow-catalog}
and~\ref{sec:workflow-cli}), the wrappers \code{FlattenJointSpaces},
\code{TeamReward} and \code{to\_parallel}
(Section~\ref{sec:workflow-wrappers}), the baseline controllers
(Section~\ref{sec:workflow-baselines}) and \code{evaluate} and \code{rollout}
(Section~\ref{sec:workflow-evaluation}).

% ---------------------------------------------------------------------------
\section{Graph utilities}\label{sec:envlib-graphs}

The module \pkg{env\_lib.utils.graphs} collects the graph computations of
networked control. A graph is a square adjacency matrix: agent $i$ interacts
with agent $j$ when entry $(i, j)$ is non-zero, and weighted matrices (such as
the line susceptances of PowerGrid) are accepted where weights make sense.
Every random topology takes a \code{numpy.random.Generator}, so graphs are
reproducible when the generator is seeded. PowerGrid builds its transmission
network with \code{make\_graph}, and the functions apply equally to the
matrices exposed by the other environments (\code{adjacency} of Consensus and
Platoon, \code{topology\_matrix} of Kuramoto).

\begin{table}[htbp]
  \centering\small
  \caption{Functions of \pkg{env\_lib.utils.graphs}.}
  \label{tab:envlib-graphs}
  \begin{tabularx}{\textwidth}{@{}>{\raggedright\arraybackslash}p{6.3cm}>{\raggedright\arraybackslash}X@{}}
    \toprule
    Function & Result \\
    \midrule
    \code{make\_graph(topology, n, *, rng=None, ...)} & symmetric boolean \code{(n, n)}
      adjacency of one of \code{TOPOLOGIES}: \code{"ring"}, \code{"line"},
      \code{"star"} (node 0 is the hub), \code{"complete"}, \code{"grid"},
      \code{"erdos\_renyi"}, \code{"small\_world"} (Watts--Strogatz),
      \code{"random\_geometric"}; random topologies are resampled until
      connected (\code{connected=True}); \code{return\_positions=True} also
      returns node coordinates \\
    \code{is\_connected(adjacency)} & whether the undirected graph is connected \\
    \code{laplacian(adjacency)} & $L = D - A$ (weighted for a weighted matrix) \\
    \code{algebraic\_connectivity(adjacency)} & Fiedler value $\lambda_2(L)$ of the
      symmetrised graph; positive exactly for a connected graph, and the rate of
      linear consensus \\
    \code{k\_hop\_adjacency(adjacency, k, *, include\_self=False)} & boolean matrix of
      the pairs within $k$ hops: the neighbourhoods of $\kappa$-hop (truncated)
      policies and critics \\
    \code{degree}, \code{neighbor\_lists}, \code{edge\_list} & weighted degrees, neighbour
      indices per node, undirected edges $(i, j)$ with $i < j$ \\
    \code{grid\_shape(n)} & rows and columns of the most square grid with $n$ nodes \\
    \code{circular\_layout(n)}, \code{spring\_layout(adjacency)} & node positions for
      plotting (a circle, or a deterministic force-directed layout in $[-1, 1]^2$) \\
    \bottomrule
  \end{tabularx}
\end{table}

\begin{pycode}
import numpy as np
import env_lib
from env_lib.utils import graphs

adj = graphs.make_graph("small_world", 12, rng=np.random.default_rng(0))
print(adj.shape, adj.dtype, graphs.is_connected(adj))  # (12, 12) bool True
print(len(graphs.edge_list(adj)), graphs.degree(adj).max())  # 24 6.0
lap = graphs.laplacian(adj)                            # L = D - A
print(round(graphs.algebraic_connectivity(adj), 3))    # 1.35
two_hop = graphs.k_hop_adjacency(adj, 2)               # pairs within two hops
print(two_hop.sum(axis=1))             # [11  9  9 11  8 11 11 10  8  9 10 11]
xy = graphs.spring_layout(adj)                         # (12, 2) in [-1, 1]

grid = env_lib.make("PowerGrid-v0").unwrapped          # the benchmark grid
print(round(graphs.algebraic_connectivity(grid.adjacency), 3))     # 0.668
print(round(graphs.algebraic_connectivity(grid.susceptance), 3))   # 0.424
\end{pycode}

The last line weights every line of the benchmark grid with its susceptance
$B_{ij}$, which gives the algebraic connectivity of the weighted Laplacian that
governs the electro-mechanical modes of the grid.

% ---------------------------------------------------------------------------
\section{Rendering}\label{sec:envlib-rendering}

The render mode is chosen at construction with the \code{render\_mode}
argument:

\begin{description}[style=nextline,leftmargin=2.8cm]
  \item[\code{None}] (default) No rendering. \code{render()} returns
    \code{None} and warns once.
  \item[\code{"rgb\_array"}] \code{render()} returns the current frame as a
    \code{(H, W, 3)} \code{uint8} array. Frames are drawn off-screen and need
    no display.
  \item[\code{"human"}] A window is opened on the first frame and updated
    automatically at the end of every \code{reset()} and \code{step()},
    following the Gymnasium convention; do not call \code{render()} yourself
    in this mode. The frame rate is capped at \code{metadata["render\_fps"]}.
  \item[\code{"ansi"}] (LineMsg and WirelessComm) \code{render()} returns a
    plain-text summary of the agents' actions and states.
\end{description}

The native vector environments accept \code{render\_mode} as well and draw
copy 0; a \code{SyncVectorEnv} returns one frame per copy.

\begin{table}[htbp]
  \centering\small
  \caption{Render modes and frame sizes (width $\times$ height in pixels).}
  \label{tab:envlib-render}
  \begin{tabularx}{\textwidth}{@{}lll>{\raggedright\arraybackslash}X@{}}
    \toprule
    Environment & Modes & Frame size & Drawing backend, \code{render\_fps} \\
    \midrule
    Kuramoto & human, rgb\_array & $1000 \times 560$ & matplotlib, 30 \\
    LineMsg & human, rgb\_array, ansi & $1000 \times 560$ & matplotlib, 10 \\
    WirelessComm & human, rgb\_array, ansi & $1000 \times 560$ & matplotlib, 10 \\
    Pistonball & human, rgb\_array & $(160 + 40n) \times 560$ & pygame, 20 \\
    Consensus / Formation & human, rgb\_array & $1000 \times 560$ & matplotlib, 20 \\
    PowerGrid & human, rgb\_array & $1000 \times 560$ & matplotlib, 20 \\
    Platoon & human, rgb\_array & $1000 \times 560$ & matplotlib, 10 \\
    AJLATT & human, rgb\_array & $1100 \times 620$ & matplotlib, 8 \\
    \bottomrule
  \end{tabularx}
\end{table}

\paragraph{Renderer design.} The matplotlib renderers are subclasses of
\code{MatplotlibRenderer} from \code{env\_lib.utils}. They create their axes and artists
once and only update the artists' data on later frames. In
\code{"rgb\_array"} mode the static part of the figure (axes, ticks, labels,
legends) is rasterised once and cached, and each frame redraws only the
registered dynamic artists (blitting), which is typically three to six times
faster than a full redraw. The renderers never modify
\code{matplotlib.rcParams}. Pistonball draws with pygame surfaces.

\paragraph{Human mode and headless machines.} Human mode needs an interactive
matplotlib backend (for example TkAgg or QtAgg) or, for Pistonball, a display
for pygame. With a non-interactive backend the environment issues a
\code{RuntimeWarning} and shows no window. On servers and in continuous
integration set
\begin{shellcode}
export MPLBACKEND=Agg SDL_VIDEODRIVER=dummy
\end{shellcode}
and use \code{render\_mode="rgb\_array"}.

\paragraph{Themes.} Two colour themes ship with the package: \code{"dark"}
(the default) and \code{"light"}, which is intended for papers and slides and
was used for all figures in this manual. Select a theme with
\code{env\_lib.utils.set\_theme}. The theme is read when a renderer is created,
that is at the first \code{render()} call of an environment, so set it before
rendering starts. Custom palettes are \code{env\_lib.utils.Theme} instances
registered with \code{register\_theme}.

\begin{pycode}
import env_lib
from env_lib.utils import set_theme

set_theme("light")          # renderers created from now on use the light theme
env = env_lib.make("Consensus-v0", render_mode="rgb_array")
obs, info = env.reset(seed=0)
for _ in range(100):
    action = env.action_space.sample()
    obs, reward, terminated, truncated, info = env.step(action)
    if terminated or truncated:
        break
frame = env.render()        # numpy array of shape (560, 1000, 3), dtype uint8
env.close()
set_theme("dark")           # back to the default theme
\end{pycode}

% ---------------------------------------------------------------------------
\section{Recording episodes}\label{sec:envlib-recording}

\code{env\_lib.utils} provides two functions to turn rollouts into animations:

% norun
\begin{pycode}
record_episode(env, policy=None, path=None, *, max_steps=None, seed=None,
               fps=None, render_every=1) -> list[numpy.ndarray]
save_animation(frames, path, fps=30.0, *, loop=0) -> pathlib.Path
\end{pycode}

\code{record\_episode} resets the environment with \code{seed}, renders the
initial state and then steps the environment with \code{policy} (a callable
mapping an observation to an action; uniformly random actions from
\code{env.action\_space} if omitted), rendering one frame every
\code{render\_every} steps. The rollout stops when any entry of
\code{terminated} or \code{truncated} is true or after \code{max\_steps}
steps. The function returns the list of frames and, when \code{path} is given,
writes them with \code{save\_animation} at \code{fps} frames per second (default
\code{env.metadata["render\_fps"]}). The environment must use
\code{render\_mode="rgb\_array"}.

\code{save\_animation} writes a sequence of equally sized RGB frames and
returns the path. The format follows the file suffix: \code{.gif} is written
with Pillow (always available); \code{.mp4}, \code{.webm}, \code{.avi},
\code{.mov} and \code{.mkv} are written with \pkg{imageio} and
\pkg{imageio-ffmpeg} (extra \code{video}); frames are cropped to even
dimensions for video codecs. \code{loop} is the GIF loop count
(0 loops forever). Missing parent directories are created.

\begin{pycode}
import env_lib
from env_lib.utils import record_episode, save_animation

env = env_lib.make("Formation-v0", render_mode="rgb_array")
policy = env_lib.baseline_policy(env)                    # Laplacian controller
frames = record_episode(env, policy, "renders/formation.gif",
                        max_steps=150, seed=0, render_every=2)
save_animation(frames[::2], "renders/formation_fast.gif", fps=20)
env.close()
\end{pycode}

The command \code{env-lib run Formation-v0 -{}-gif renders/formation.gif}
records the same kind of episode without writing code
(Section~\ref{sec:workflow-cli}).

% ---------------------------------------------------------------------------
\section{Registered environments}\label{sec:envlib-ids}

Table~\ref{tab:envlib-ids} lists all 18 ids returned by
\code{env\_lib.list\_envs()} with the constructor arguments each registration
sets; all other arguments keep their defaults. The last column shows how
\code{env\_lib.make\_vec} vectorises the id: with a native implementation or
with \code{SyncVectorEnv}. The version suffixes follow the historical ids of
the project and denote preset configurations, not revisions:
\envid{WirelessComm-v1}, for example, is the $4 \times 4$ grid and not a newer
version of \envid{WirelessComm-v0}. Gymnasium may therefore emit an ``out of
date'' \code{DeprecationWarning} when a \code{-v0} id is created with
\code{gymnasium.make}; the warning is harmless, and \code{env\_lib.make}
suppresses it.

\begingroup\small
\begin{longtable}{@{}>{\raggedright\arraybackslash}p{0.43\textwidth}>{\raggedright\arraybackslash}p{0.43\textwidth}>{\raggedright\arraybackslash}p{0.08\textwidth}@{}}
  \caption{Registered Gymnasium ids and their configuration.}\label{tab:envlib-ids}\\
  \toprule
  Id & Class and configuration & Vector \\
  \midrule
  \endfirsthead
  \toprule
  Id & Class and configuration & Vector \\
  \midrule
  \endhead
  \bottomrule
  \endfoot
  \envid{KuramotoOscillator-v0} & \code{KuramotoOscillatorEnv}; defaults (10 oscillators, dynamic coupling) & native \\
  \envid{KuramotoOscillator-v1} & \code{KuramotoOscillatorEnv}; \code{n\_oscillators=5} & native \\
  \envid{KuramotoOscillator-Constant-v0} & \code{KuramotoOscillatorEnv}; \code{n\_oscillators=6}, \code{coupling\_mode="constant"} & native \\
  \envid{KuramotoOscillator-\allowbreak FreqSync-Constant-v0} & \code{KuramotoOscillatorEnv}; \code{n\_oscillators=6}, \code{coupling\_mode="constant"}, \code{reward\_type=}\allowbreak\code{"frequency\_synchronization"}, \code{target\_frequency=2.0} & native \\
  \envid{KuramotoOscillatorTorch-v0} & \code{KuramotoOscillatorEnvTorch}; defaults (10 oscillators, one system) & sync \\
  \envid{KuramotoOscillatorTorch-v1} & \code{KuramotoOscillatorEnvTorch}; \code{n\_oscillators=8}, \code{n\_agents=4}, \code{device="cpu"} & sync \\
  \envid{KuramotoOscillatorTorch-v2} & \code{KuramotoOscillatorEnvTorch}; \code{n\_oscillators=10}, \code{n\_agents=1}, \code{device="auto"}, \code{integration\_method="rk4"} & sync \\
  \envid{KuramotoOscillatorTorch-\allowbreak Constant-v0} & \code{KuramotoOscillatorEnvTorch}; \code{n\_oscillators=6}, \code{coupling\_mode="constant"} & sync \\
  \envid{KuramotoOscillatorTorch-\allowbreak FreqSync-Constant-v0} & \code{KuramotoOscillatorEnvTorch}; \code{n\_oscillators=6}, \code{coupling\_mode="constant"}, \code{reward\_type=}\allowbreak\code{"frequency\_synchronization"}, \code{target\_frequency=2.0} & sync \\
  \envid{LineMsg-v0} & \code{LineMsgEnv}; defaults (10 agents) & sync \\
  \envid{WirelessComm-v0} & \code{WirelessCommEnv}; defaults ($6 \times 6$ grid) & sync \\
  \envid{WirelessComm-v1} & \code{WirelessCommEnv}; \code{grid\_x=4}, \code{grid\_y=4} & sync \\
  \envid{Pistonball-v0} & \code{PistonballEnv}; defaults (20 pistons); needs \pkg{pymunk} (and \pkg{pygame} to render) & sync \\
  \envid{Consensus-v0} & \code{ConsensusEnv}; defaults (8 agents, rendezvous on a ring) & native \\
  \envid{Formation-v0} & \code{ConsensusEnv}; \code{task="formation"} (circle formation) & native \\
  \envid{PowerGrid-v0} & \code{PowerGridEnv}; defaults (16 buses, small-world network fixed by \code{network\_seed=0}) & native \\
  \envid{Platoon-v0} & \code{PlatoonEnv}; defaults (8 followers, predecessor following, mixed leader scenario) & native \\
  \envid{AJLATT-v0} & \code{AJLATTEnv}; defaults (4 robots on \code{obstacles04}); passive environment checker disabled because rewards and \code{terminated} are per-robot arrays & sync \\
\end{longtable}
\endgroup

\endgroup
